% !TeX root = ../main.tex
\\

\textbf{Тестовое воздействие сигналом с линейно возрастающей амплитудой (КИХ-дециматор)}

Для проверки линейности тракта с единым КИХ-дециматором и визуализации его амплитудно-частотной характеристики (АЧХ) применялся монохроматический сигнал с линейно нарастающей огибающей (LA). Частота тестового сигнала была смещена на 1 МГц относительно частоты гетеродина, чтобы результат демодуляции попал в рабочую полосу пропускания фильтра. Максимальная амплитуда ограничивалась значением 255 LSB для соответствия динамическому диапазону АЦП.

Результаты моделирования во временной и частотной областях представлены на рисунках 14 — 17.

\begin{figure}[H]
    \centering
    \includegraphics[width=0.8\textwidth]{P5_Signal_LA_on_IF.png}
    \caption{Временная диаграмма входного сигнала с линейно возрастающей амплитудой на ПЧ}
    \label{fig:p7_la_sig_in}
\end{figure}

\begin{figure}[H]
    \centering
    \includegraphics[width=0.8\textwidth]{P5_Signal_LA_out.png}
    \caption{Временная диаграмма демодулированного и отфильтрованного сигнала (I и Q каналы)}
    \label{fig:p7_la_sig_out}
\end{figure}

\begin{figure}[H]
    \centering
    \includegraphics[width=0.8\textwidth]{P5_Spectrum_LA_on_IF.png}
    \caption{Спектр входного сигнала на ПЧ}
    \label{fig:p7_la_spec_in}
\end{figure}

\begin{figure}[H]
    \centering
    \includegraphics[width=0.8\textwidth]{P5_Spectrum_LA_out.png}
    \caption{Спектр на выходе КИХ-дециматора (визуализация АЧХ фильтра)}
    \label{fig:p7_la_spec_out}
\end{figure}

\textbf{Выводы по результатам моделирования:}
Временная диаграмма сигнала на выходе КИХ-дециматора (рис. 15) демонстрирует идеальное сохранение линейного закона нарастания амплитуды для синфазного и квадратурного каналов. Отсутствие искажений огибающей и переполнений свидетельствует о корректном расчете коэффициентов фильтра и достаточной разрядности аккумуляторов при выполнении операции свертки.

Спектр выходного сигнала (рис. 17) наглядно отражает профиль АЧХ спроектированного КИХ-фильтра. В отличие от каскада Хогенауэра, здесь наблюдается абсолютно плоская вершина в полосе пропускания без необходимости применения дополнительных компенсаторов. Полезный сигнал в виде острого пика на частоте 1 МГц передан без амплитудных потерь, а широкополосный шум за пределами частоты среза резко подавлен, формируя характерную прямоугольную спектральную маску с высокой крутизной скатов.